% !TEX root = ../print.tex
% Draft \S6 --- Scale covariance: the similarity form
%
% Source: draft/spatial-hemigroup-scale-space.md, \S6 (lines 280-364).
%
% Five nodes. lem:covariance-fourier and the bulk of lem:action-rigidity and
% prop:canonical-gauge port verbatim from the causal chapter. Three things are new and are
% exactly the three places the causal proof used its vanishing lemma:
%
%   lem:no-lattice      -- the a priori exclusion of lattice kernels, from covariance at a
%                          single non-integer ratio and the chain structure of the zero sets;
%   lem:action-rigidity(3) -- continuity of the action, read off the smoothed transmittance
%                          rather than off the accumulated exponent at one frequency;
%   lem:dilation-atom   -- the direction of the orbit coordinate, where the causal argument
%                          used monotonicity of Bernstein functions, unavailable here before
%                          the classification (it is recovered in prop:strict-positivity(1)).
%
% CHECK (draft, checks for \S6, item 3): the iteration "a function invariant under one
% dilation different from 1 and continuous at 0 with value 0 there vanishes identically" is
% used three times below (lem:action-rigidity(4), the injectivity of c, and the uniqueness
% clause of thm:main-characterization). It is stated once, as lem:dilation-invariance.

\chapter{Scale covariance: the similarity form}
\label{ch:covariance}

Covariance turns the free relabelling $S_\lambda$ of (A8) into a coordinate: the action is
rigid, its orbit through $1$ parametrizes the scale axis, and in the resulting gauge the
accumulated exponent is a dilate of a single function. In the reading of
\ref{rem:modulation-transfer}, covariance makes the family of transfer functions
\emph{self-similar} --- descending the scale axis is indistinguishable from probing at a finer
wavenumber --- and the rigidity lemma says that no scale is distinguished by the action.

% draft: Lemma 6.1'
% twin: hcs lem:covariance-laplace (Hemigroup.CascadeCore.covariance_laplace) --- verbatim.
% The Lean statement (Skeleton.covariance_fourier) takes the setting the proviso in clause (3)
% describes -- (A6)--(A7) present, so the exponents exist -- and is therefore unconditional
% there. The (A1)--(A5)-only form is not stated in Lean because nothing consumes it.
\begin{lemma}[Fourier form of covariance]\label{lem:covariance-fourier}
\uses{def:cascade-family,lem:convolution-representation,lem:nonvanishing}
\lean{SpatialLine.covariance_fourier, SpatialLine.covariance_similarity}\leanok
\notes{covariance-strata}
Under (A1)--(A5) the following are equivalent: (1) axiom (A8); (2)
$\Dil{\lambda}\mu_{s,t} = \mu_{S_\lambda s,\,S_\lambda t}$ for all $0 \le s \le t$ and
$\lambda > 0$; (3) $g_{S_\lambda s,\,S_\lambda t}(\omega) = g_{s,t}(\lambda\omega)$ for all
$0 \le s \le t$, $\lambda > 0$ and $\omega \in \RR$, whenever the exponents exist. In
particular, under (A8) together with (A6)--(A7),
\begin{equation}\label{eq:similarity}
  G(S_\lambda t,\,\omega) = G(t,\,\lambda\omega)
  \qquad \text{for all } t \ge 0,\ \lambda > 0,\ \omega \in \RR .
\end{equation}
\statusT\quad\emph{(The proviso in clause (3) is there because (A1)--(A5) alone do not give
\ref{lem:nonvanishing}; under (A6)--(A7) it is vacuous, and every later use is in that
setting. \textbf{Machine-checked}, \texttt{\#print axioms} reducing to Lean core, in the
setting the proviso describes: the formal statement carries (A6)--(A7) and a kernel family,
which is what makes clause (3) unconditional, and the equivalence is stated on the
intertwining alone, the three conditions on $S_\lambda$ being hypotheses shared by all three
clauses.)}
\end{lemma}

\begin{proof}
\leanok
The compatibility $\Dil{\lambda}(\mu * f) = (\Dil{\lambda}\mu) * (\Dil{\lambda}f)$ turns (A8)
into the statement that $\Dil{\lambda}\mu_{s,t}$ and $\mu_{S_\lambda s,S_\lambda t}$ represent
the same operator, and the uniqueness clause of \ref{lem:convolution-representation} makes
that an identity of measures; this is (1) if and only if (2). For (2) if and only if (3),
$\widehat{\Dil{\lambda}\mu}(\omega) = \hat\mu(\lambda\omega)$, so taking $-\log$
(\ref{lem:nonvanishing}) gives $g_{S_\lambda s,S_\lambda t}(\omega) = g_{s,t}(\lambda\omega)$.
Equation \ref{eq:similarity} is the case $s = 0$ of that identity together with
$S_\lambda 0 = 0$, which holds because an increasing bijection of $[0,\infty)$ fixes $0$.
\end{proof}

The first thing covariance buys on the line is the exclusion of the lattice laws that
\ref{lem:lattice-zero} left open.

% draft: Lemma 6.0' (stated in \S6, after the Fourier form of covariance it uses)
% NEW: no causal counterpart; on the half-line the corresponding statement is the vanishing
% lemma, hcs lem:vanishing, which is a property of the function class and needs no axiom.
\begin{lemma}[the kernels from the origin are not lattice laws]\label{lem:no-lattice}
\uses{lem:covariance-fourier,lem:lattice-zero,lem:additivity,lem:nonvanishing}
\lean{SpatialLine.no_lattice}\leanok
Under (A1)--(A8) and (ND), for every $t > 0$ the kernel $\mu_{0,t}$ is carried by no lattice
$p\ZZ$ with $p > 0$; equivalently $G(t,\omega) > 0$ for every $\omega \ne 0$.
\statusT\quad\emph{(Covariance at a single non-integer ratio, and nothing else beyond the
cascade. The argument does not extend to the increments $\mu_{s,t}$ with $s > 0$, because
$S_\lambda$ moves both endpoints and the increments do not form a chain; strict positivity for
every increment is recovered after the classification, in
\ref{prop:strict-positivity}(2). That covariance is the operative axiom is shown by
\ref{rem:lattice-family}, a family satisfying (A1)--(A7) and (ND) all of whose kernels are
lattice laws. \textbf{Machine-checked}, \texttt{\#print axioms} reducing to Lean core. The
formal comparison of the two chain members is made element-wise rather than by classifying the
inclusions: whichever way the chain runs, one generator lands in the other lattice and produces
$3/2 = n$ or $2/3 = n$ in $\ZZ$.)}
\end{lemma}

\begin{proof}
\leanok
Write $N_t := N(\mu_{0,t})$ in the notation of \ref{lem:lattice-zero}, a closed subgroup of
$\RR$. For $s \le t$, $g_{0,t} = g_{0,s} + g_{s,t}$ with both terms nonnegative
(\ref{lem:additivity}, \ref{lem:nonvanishing}), so $N_t \subseteq N_s$: the family
$(N_t)_{t>0}$ is a chain under inclusion, and $N_t \ne \RR$ for $t > 0$ by (ND) and
\ref{lem:lattice-zero}. By \ref{lem:covariance-fourier},
$\hat\mu_{0,S_\lambda t}(\omega) = \hat\mu_{0,t}(\lambda\omega)$, so
$N_{S_\lambda t} = \lambda^{-1}N_t$, and $S_\lambda t > 0$. Suppose $N_t = c\ZZ$ for some
$t > 0$ and $c > 0$. Then for every $\lambda > 0$ the two chain members $\lambda^{-1}c\ZZ$ and
$c\ZZ$ are comparable under inclusion, which happens if and only if $\lambda$ or
$\lambda^{-1}$ is a positive integer; $\lambda = 3/2$ is neither. So $N_t = \{0\}$, which is
the statement.
\end{proof}

% additive: the iteration the draft performs three times (its Lemma 6.2'(4), the injectivity
% of c in Proposition 6.3', and the uniqueness clause of Theorem 7.3'), stated once as the
% draft's own check for \S6 (item 3) suggests.
\begin{lemma}[a function fixed by one dilation vanishes]\label{lem:dilation-invariance}
\lean{SpatialLine.dilation_invariance}\leanok
Let $h : \RR \to \RR$ be continuous at $0$ with $h(0) = 0$, and suppose
$h(\omega) = h(\varkappa\omega)$ for every $\omega$ and one fixed $\varkappa > 0$ with
$\varkappa \ne 1$. Then $h \equiv 0$.
\statusT\quad\emph{(Three lines, stated once so that the three uses cite it instead of
repeating it. \textbf{Machine-checked}, \texttt{\#print axioms} reducing to Lean core. The Lean
proof iterates over $n \in \NN$ and reduces the expanding case $\varkappa > 1$ to the
contracting one by replacing $\varkappa$ with $\varkappa^{-1}$, which is the same argument
written without a sign choice.)}
\end{lemma}

\begin{proof}
\leanok
Iterating, $h(\omega) = h(\varkappa^n\omega)$ for every $n \in \ZZ$. Choose the sign of $n$ so
that $\varkappa^n \to 0$; then $\varkappa^n\omega \to 0$ and continuity at $0$ gives
$h(\omega) = h(0) = 0$.
\end{proof}

Nothing in (A8) says that the relabellings $S_\lambda$ are unique, that they compose, or that
they move every point. The next lemma says all three. The causal engine was the injectivity of
the accumulated exponent in scale at one fixed value of the transform variable, which
\ref{cor:monotonicity} does not deliver here. Two weaker facts replace it and suffice: the
accumulated exponent is injective in scale \emph{as a function of $\omega$}, which is (ND)
restated, and the smoothed transmittance is injective in scale at a single number.

% draft: Lemma 6.2'
% twin: hcs lem:action-rigidity (Hemigroup.CascadeCore.action_rigidity); clauses (2) and (4)
% port verbatim, clauses (1) and (3) have new proofs.
\begin{lemma}[rigidity of the action]\label{lem:action-rigidity}
\uses{lem:covariance-fourier,cor:smoothed-transmittance,lem:convolution-representation,lem:additivity,lem:dilation-invariance}
\lean{SpatialLine.action_rigidity_injective, SpatialLine.action_rigidity_group,
SpatialLine.action_rigidity_continuous, SpatialLine.action_rigidity_no_fixed_point}\leanok
\notes{covariance-strata}
Assume (A1)--(A8) and (ND). Then:
\begin{enumerate}
\item If $G(t_1,\cdot) = G(t_2,\cdot)$ as functions on $\RR$, then $t_1 = t_2$; consequently,
for each $\lambda$, $S_\lambda$ is uniquely determined by \ref{eq:similarity}.
\item $S_\lambda S_\kappa = S_{\lambda\kappa}$ for all $\lambda,\kappa > 0$, and
$S_1 = \mathrm{Id}$.
\item $\lambda \mapsto S_\lambda t$ is continuous for each $t > 0$.
\item The action has no fixed point in $(0,\infty)$, in the strong form: if
$S_\lambda t^* = t^*$ for one $\lambda \ne 1$, then $t^* = 0$.
\end{enumerate}
\statusT\quad\emph{(Three of the four clauses are the function-injectivity of (1) and nothing
else. Clause (3) is the one that needs more, and what it needs is a single strictly monotone
number attached to each scale, which is what $\Theta$ is. Clause (4) is where (ND) bites in
this chapter, exactly as in the causal argument: a fixed point of the action would be a scale
at which the family is the identity, an idle stretch, and the orbit below could stall there.
The strong form, one $\lambda$ rather than all, costs nothing and is used in
\ref{prop:canonical-gauge}. \textbf{Machine-checked}, \texttt{\#print axioms} reducing to Lean
core, one declaration per clause. One remark on clause (3): the proof below inverts $\Theta$,
and the formalisation does not, because it does not have to. The order characterisation of a
limit in $\RR$ splits the claim into the two one-sided bounds $a < S_\lambda t$ and
$S_\lambda t < b$, and each follows from the strict monotonicity of $\Theta$ by contraposition
against a bound the continuity of $\theta_t$ already supplies; no inverse function is
constructed.)}
\end{lemma}

\begin{proof}
\leanok
\emph{(1).} If $t_1 < t_2$ then $G(t_2,\cdot) - G(t_1,\cdot) = g_{t_1,t_2}$ by
\ref{lem:additivity}; if this vanishes identically then $\hat\mu_{t_1,t_2} \equiv 1$, so
$\mu_{t_1,t_2} = \delta_0$ and $\Phi_{t_1,t_2} = \mathrm{Id}$ by
\ref{lem:convolution-representation}, against (ND). If $S$ and $S'$ both satisfy
\ref{eq:similarity} then $G(S_\lambda t,\cdot) = G(t,\lambda\,\cdot) = G(S'_\lambda t,\cdot)$,
so $S_\lambda t = S'_\lambda t$.

\emph{(2).} Applying \ref{eq:similarity} twice,
$G(S_\lambda S_\kappa t,\omega) = G(S_\kappa t,\lambda\omega) = G(t,\lambda\kappa\omega)
= G(S_{\lambda\kappa}t,\omega)$ for every $\omega$, and (1) gives
$S_\lambda S_\kappa t = S_{\lambda\kappa}t$. Taking $\lambda = 1$ in \ref{eq:similarity} gives
$G(S_1t,\cdot) = G(t,\cdot)$, whence $S_1 = \mathrm{Id}$.

\emph{(3).} Fix $t > 0$ and let $\rho$ be the standard Gaussian density. By
\ref{eq:similarity} and the definition of $\Theta$ in \ref{cor:smoothed-transmittance},
\[
  \Theta(S_\lambda t)
  = \int_\RR \hat\mu_{0,S_\lambda t}(\omega)\,\rho(\omega)\,d\omega
  = \int_\RR \hat\mu_{0,t}(\lambda\omega)\,\rho(\omega)\,d\omega
  = \int_\RR \hat\mu_{0,t}(u)\,\rho(u/\lambda)\,\lambda^{-1}\,du \;=:\; \theta_t(\lambda).
\]
The last integrand is continuous in $\lambda$ and, for $\lambda$ in a compact subinterval
$[\lambda_0/2, 2\lambda_0]$ of $(0,\infty)$, is dominated by
$\tfrac{2}{\lambda_0}(2\pi)^{-1/2}e^{-u^2/(8\lambda_0^2)}$, which is integrable; so $\theta_t$
is continuous by dominated convergence. Under (ND), $\Theta$ is continuous and strictly
decreasing on $[0,\infty)$ (\ref{cor:smoothed-transmittance}), hence has a continuous inverse
on its image, and $\theta_t(\lambda) = \Theta(S_\lambda t)$ lies in that image; so
$S_\lambda t = \Theta^{-1}(\theta_t(\lambda))$ is continuous in $\lambda$.

\emph{(4).} Suppose $S_\lambda t^* = t^*$ for some $\lambda \ne 1$. Then \ref{eq:similarity}
gives $G(t^*,\omega) = G(t^*,\lambda\omega)$ for every $\omega$, and $G(t^*,\cdot)$ is
continuous with $G(t^*,0) = 0$ (\ref{lem:additivity}), so $G(t^*,\cdot) \equiv 0$ by
\ref{lem:dilation-invariance}; that is, $\mu_{0,t^*} = \delta_0$ and
$\Phi_{0,t^*} = \mathrm{Id}$, which by (ND) forces $t^* = 0$.
\end{proof}

The direction of the orbit coordinate needs one more elementary fact, the Fourier form of "a
law whose dilates converge to a point mass is a point mass".

% draft: Lemma 6.2''
% NEW: the replacement for the causal argument, which oriented the orbit coordinate by the
% monotonicity of Bernstein functions. Monotonicity of F in omega is not available here before
% the classification; it is recovered from it, in prop:strict-positivity(1).
\begin{lemma}[dilation to infinity keeps only the atom at the origin]\label{lem:dilation-atom}
\lean{SpatialLine.dilation_atom}\leanok
Let $\mu$ be a probability measure on $\RR$ and let $\lambda_n \to \infty$. If
$\hat\mu(\lambda_n\omega) \to 1$ for every $\omega \in \RR$, then $\mu = \delta_0$.
\statusT\quad\emph{(The Fourier form of "$\Dil{\lambda_n}\mu \to \delta_0$ weakly implies
$\mu = \delta_0$", proved so as to need no weak convergence at all: two applications of
dominated convergence with the constant bound $1$. \textbf{Machine-checked},
\texttt{\#print axioms} reducing to Lean core. The formal proof reaches the same conclusion by
a different route, and the reason is worth recording: the truncation inequality
$\mu\{|x| > r\} \le \tfrac r2\,\bigl|\int_{-2/r}^{2/r}(1 - \hat\mu)\bigr|$ is already in
Mathlib, so applying it to the dilate of $\mu$ by $\lambda_n$ at the fixed radius $1$ leaves
one dominated convergence over a fixed interval and needs neither the Gaussian nor Fubini. The
proof written here is the proof of record; the substitution is an implementation choice.)}
\end{lemma}

\begin{proof}
\leanok
With $\rho$ the standard Gaussian density, Fubini gives
$\int \hat\mu(\lambda_n\omega)\rho(\omega)\,d\omega = \int e^{-\lambda_n^2x^2/2}\,\mu(dx)$. The
left side tends to $\int\rho = 1$ by dominated convergence, since $|\hat\mu| \le 1$; the right
side tends to $\mu(\{0\})$ by dominated convergence, the integrand being bounded by $1$ and
tending pointwise to the indicator of $\{0\}$. So $\mu(\{0\}) = 1$.
\end{proof}

The coordinate itself is the orbit. Fix the point $1$ on the scale axis and follow it under the
action: $\lambda \mapsto S_\lambda 1$ is continuous, injective, and, having no fixed point to
stall at, sweeps the whole axis; the one thing to settle is which way it sweeps.

% draft: Proposition 6.3'
% twin: hcs prop:canonical-gauge (Hemigroup.CascadeCore.canonical_gauge); the direction clause
% has a new proof.
% CHANGED (proof of record 2026-09-09): the surjectivity step now follows the machine-checked
% route, per CLAUDE.md's convention. The old argument took the limits of c at 0 and at infinity
% and asserted that each is fixed by every S_kappa "S_kappa being continuous" -- continuity of
% S_kappa in the SCALE variable, which this chapter never proves. The new one fixes u > 0 and
% argues by contradiction from the infimum and the supremum of the orbit, using only that
% S_kappa is nondecreasing: c(lambda) = S_kappa(c(lambda/kappa)) bounds S_kappa m below m and
% S_kappa M above M, and reading each at kappa and kappa^{-1} makes it a fixed point, which
% lem:action-rigidity(4) forbids; the intermediate value theorem then closes it. This is
% SpatialLine.canonical_gauge_orbit's proof. Statement unchanged; \uses unchanged.
% CHANGED (proof of record 2026-09-09): the DIRECTION clause now follows the machine-checked
% route too, closing the divergence the previous marker left open (wave 2's merge). The old
% argument proved injectivity from lem:dilation-invariance, deduced strict monotonicity from
% "continuous and injective on a connected set", and excluded the decreasing case by an appeal
% to continuity of S_lambda in the SCALE variable. The new one proves the single statement
% SpatialLine.lt_S_of_one_lt -- for kappa > 1 and u > 0, u < S_kappa u -- by running the old
% decreasing-case argument on the orbit of u, with the limit shown to be a fixed point by
% monotonicity alone; strict monotonicity of lambda |-> S_lambda t for EVERY t > 0, injectivity
% included, is then one line of the group law (SpatialLine.S_strictMonoOn_ratio). Statement
% unchanged. \uses LOSES lem:dilation-invariance, which the rewritten proof no longer cites;
% lem:dilation-atom and lem:nonvanishing are still spent, in the expanding step.
\begin{proposition}[orbit coordinate; the gauge]\label{prop:canonical-gauge}
\uses{lem:action-rigidity,lem:dilation-atom,lem:covariance-fourier,thm:increments-levy,lem:nonvanishing}
\lean{SpatialLine.canonical_gauge_orbit, SpatialLine.canonical_gauge}\leanok
\notes{covariance-strata}
Assume (A1)--(A8) and (ND). Then $c(\lambda) := S_\lambda 1$ is a continuous strictly
increasing bijection of $(0,\infty)$ onto $(0,\infty)$ with $c(1) = 1$. Consequently
\[
  \chi := c^{-1} : (0,\infty) \to (0,\infty), \qquad \chi(0) := 0,
\]
is an increasing bijection of $[0,\infty)$, the \emph{canonical gauge}, satisfying
$\chi(S_\lambda t) = \lambda\,\chi(t)$; the maps $\lambda \mapsto S_\lambda t$ are strictly
increasing for every $t > 0$; and
\begin{equation}\label{eq:gauge}
  G(t,\omega) = F\bigl(\chi(t)\,\omega\bigr), \qquad F := G(1,\cdot) \in \LEs,\ F \not\equiv 0,
\end{equation}
with $F$ even and continuous, $F(0) = 0$, and with the pair $(a,\nu)$ of
\ref{eq:levy-khintchine}.
\statusT\quad\emph{(The gauge is recovered, not posited; it is the semigroup literature's free
rescaling function, and the hemigroup case differs from the semigroup case only by the
generality of the reparametrization and the freedom in $F$. In the canonical gauge every
increment reads $g_{s,t}(\omega) = F(t\omega) - F(s\omega)$, which is the form
Chapter~\ref{ch:characterization} classifies. \textbf{Machine-checked}, \texttt{\#print axioms}
reducing to Lean core, in two declarations: the orbit coordinate --- that $c$ is a continuous
strictly increasing bijection of $(0,\infty)$ with $c(1) = 1$ --- and the gauge with the
similarity form, which also states the normalization $\chi(1) = 1$. That conjunct is immediate
from $c(1) = 1$, and it is stated because \ref{thm:main-characterization} has no other route to
it: it is the normalization that theorem's uniqueness clause quantifies over, and the necessity
direction reaches it only through this gauge.
The clause $F \in \LEs$ is \ref{thm:increments-levy} at $(0,1)$, whose own
footprint is Lean core, so no interface is spent here. Two riders are read rather than restated
in the formalisation, so that a reader of the declaration list does not look for names that do
not exist: that $F$ is even and continuous with $F(0) = 0$ is
\ref{lem:nonvanishing}'s own declaration \texttt{SpatialLine.nonvanishing\_exponent} applied at
$(0,1)$, and that $F$ carries a pair $(a,\nu)$ of \ref{eq:levy-khintchine} is the definition of
$\LEs$ unfolded. The surjectivity step below is the machine-checked proof, written out. It replaces
an argument that read the limits of $c$ at $0$ and at $\infty$ off the continuity of $S_\lambda$
in the \emph{scale} variable --- a fact this chapter does not prove --- by one that uses the
monotonicity of $S_\kappa$ alone: the extrema of the orbit satisfy $S_\kappa m \le m$ and
$M \le S_\kappa M$ for every $\kappa$, directly from
$c(\lambda) = S_\kappa\bigl(c(\lambda/\kappa)\bigr)$, and reading that at $\kappa$ and at
$\kappa^{-1}$ makes each a fixed point. The direction clause below is the machine-checked proof
too, and with it the last divergence closes. The formalisation obtains the direction and the
strict monotonicity of $\lambda \mapsto S_\lambda t$, for every $t > 0$ at once, from a single
statement --- for $\varkappa > 1$ and $u > 0$, $u < S_\varkappa u$, which is the old
decreasing-case argument run on the orbit of $u$ --- so the printed proof no longer appeals to
"continuous and injective on a connected set, hence strictly monotone" nor, inside the excluded
decreasing case, to the continuity of $S_\lambda$ in the scale variable, which this chapter does
not prove. Injectivity is contained in the strict monotonicity and is no longer argued
separately, so \ref{lem:dilation-invariance} leaves this node's uses: it was cited by the
injectivity paragraph alone.)}
\end{proposition}

\begin{proof}
\leanok
\emph{$c$ is continuous} by \ref{lem:action-rigidity}(3), and $c(1) = S_11 = 1$ by
\ref{lem:action-rigidity}(2).

\emph{The action expands above $1$: for every $\varkappa > 1$ and every $u > 0$,
$u < S_\varkappa u$.} Suppose not. Equality is impossible, because
\ref{lem:action-rigidity}(4) leaves $S_\varkappa$ no fixed point but $0$ and $u > 0$; so
$S_\varkappa u < u$ for some such pair. Write $w_n := S_{\varkappa^n}u$ for the orbit of $u$.
The group law gives $w_{n+1} = S_\varkappa w_n$, and $S_\varkappa$ is increasing, so
$w_1 < w_0 = u$ propagates by induction to $w_{n+1} < w_n$ for every $n$. The orbit is
nonnegative, hence decreases to $L := \inf_n w_n \ge 0$, and that infimum is a fixed point of
$S_\varkappa$ by monotonicity alone: $S_\varkappa L \le S_\varkappa w_n = w_{n+1}$ for every
$n$, so $S_\varkappa L \le L$; and $S_{\varkappa^{-1}}L \le S_{\varkappa^{-1}}w_{n+1} = w_n$
for every $n$, so $S_{\varkappa^{-1}}L \le L$ and therefore
$L = S_\varkappa\bigl(S_{\varkappa^{-1}}L\bigr) \le S_\varkappa L$. Thus $S_\varkappa L = L$,
and \ref{lem:action-rigidity}(4) gives $L = 0$: the orbit of $u$ decreases to $0$. By
\ref{eq:similarity}, $\hat\mu_{0,u}(\varkappa^n\omega) = \hat\mu_{0,w_n}(\omega)$, which tends
to $\hat\mu_{0,0}(\omega) = 1$ for every $\omega$, by continuity of
$\hat\mu_{0,\cdot}(\omega)$ in scale (\ref{lem:nonvanishing}). Then \ref{lem:dilation-atom}
gives $\mu_{0,u} = \delta_0$, against (ND).

\emph{$c$ is strictly increasing}, and so is $\lambda \mapsto S_\lambda t$ for every $t > 0$.
Let $0 < a < b$ and $t > 0$. The group law reads $S_b t = S_{b/a}(S_a t)$ with $b/a > 1$ and
$S_a t > 0$, so the previous paragraph gives $S_a t < S_b t$. Taking $t = 1$ is the assertion
for $c$; injectivity is contained in it, and nothing here needs $c$ to be continuous.

\emph{$c$ is onto $(0,\infty)$.} Fix $u > 0$. Since $c$ is continuous and strictly increasing on
$(0,\infty)$, it is enough to exhibit one value of $c$ below $u$ and one above it, the
intermediate value theorem then supplying a $\lambda$ with $c(\lambda) = u$. Both are obtained
from the same two lines, and neither uses any continuity of $S_\kappa$ in the scale variable.

Suppose $c \ge u$ on $(0,\infty)$, and let $m := \inf_{\lambda > 0} c(\lambda)$, so
$m \ge u > 0$. For every $\kappa > 0$ and every $\lambda > 0$ the group law gives
$c(\lambda) = S_\kappa\bigl(c(\lambda/\kappa)\bigr) \ge S_\kappa m$, because $S_\kappa$ is
nondecreasing on $[0,\infty)$ and $c(\lambda/\kappa) \ge m$; taking the infimum over $\lambda$,
$S_\kappa m \le m$ for every $\kappa > 0$. Reading that at $\kappa^{-1}$ and applying the
nondecreasing $S_\kappa$ to it, $m = S_\kappa\bigl(S_{\kappa^{-1}}m\bigr) \le S_\kappa m$, so
$S_\kappa m = m$. Taking $\kappa = 2$, \ref{lem:action-rigidity}(4) gives $m = 0$, contradicting
$m \ge u > 0$.

Suppose instead $c \le u$ on $(0,\infty)$, and let $M := \sup_{\lambda > 0} c(\lambda)$, which
is finite and at least $c(1) = 1$. The same computation with the inequalities reversed gives
$c(\lambda) = S_\kappa\bigl(c(\lambda/\kappa)\bigr) \le S_\kappa M$ and hence
$M \le S_\kappa M$ for every $\kappa > 0$; reading that at $\kappa^{-1}$ and applying $S_\kappa$,
$S_\kappa M \le S_\kappa\bigl(S_{\kappa^{-1}}M\bigr) = M$, so $S_\kappa M = M$ and
\ref{lem:action-rigidity}(4) gives $M = 0$, contradicting $M \ge 1$.

\emph{The gauge relation.} By the group law,
$S_\lambda t = S_\lambda S_{\chi(t)}1 = S_{\lambda\chi(t)}1 = c(\lambda\chi(t))$, so
$\chi(S_\lambda t) = \lambda\chi(t)$, and $\lambda \mapsto S_\lambda t$ is strictly increasing
because $c$ is.

\emph{The similarity form.} For $t > 0$, writing $t = S_{\chi(t)}1$ and using
\ref{eq:similarity}, $G(t,\omega) = G(S_{\chi(t)}1,\omega) = G(1,\chi(t)\omega) =
F(\chi(t)\omega)$. Here $F = g_{0,1} \in \LEs$ by \ref{thm:increments-levy}, even and
continuous with $F(0) = 0$ by \ref{lem:nonvanishing}, and $F \not\equiv 0$ by (ND). The case
$t = 0$ is $G(0,\cdot) = 0 = F(0)$.
\end{proof}

\begin{remark}[gauge freedom]\label{rem:gauge-freedom}
In the hemigroup setting the scale coordinate carries no intrinsic parametrization: replacing
$t$ by any increasing bijection of $[0,\infty)$ yields an equivalent family.
\ref{prop:canonical-gauge} fixes the \emph{canonical gauge} $\tilde t := \chi(t)$, in which
$S_\lambda$ acts by multiplication and $G(t,\omega) = F(\tilde t\,\omega)$. In it scale carries
the dimension of length. Since $\omega$ is dual to position, $\tilde t\omega$ must be
dimensionless for $F(\tilde t\omega)$ to make sense, so $t$ in this article is a length, the
standard deviation of the Gaussian member, where the scale-space literature's $t$ is the
variance. The other homogeneous gauges are $c_0\,\tilde t^{\,\beta}$ with $c_0, \beta > 0$. The
one that recurs is the \emph{variance gauge} $\tilde t^{\,2}$, in which the Gaussian member is a
one-parameter semigroup and the diffusion equation has its textbook form. It is named at each
use and never written $\tilde t$.
\end{remark}

For the Gaussian member, $G(t,\omega) = \tfrac12t^2\omega^2$ in the coordinate of
Chapter~\ref{ch:axioms}, so $F(\omega) = \tfrac12\omega^2$ and that coordinate already is the
canonical gauge: it is the standard deviation. In the semigroup case the straightening is the
classical reparametrization: in the semigroup's own additive parameter $v$ the accumulated
exponent is $v\,|\omega|^\alpha$, the action is $S_\lambda v = \lambda^\alpha v$, and the gauge
is $\chi(v) = v^{1/\alpha}$, which is exactly the rescaling function of
\sscite{pauwels1995extended}; their functional equations for the exponent are the Cauchy step
of \ref{cor:semigroup-case}.
